\documentclass[10pt,twocolumn]{article}

% ── packages ──────────────────────────────────────────────────────────────────
\usepackage[utf8]{inputenc}
\usepackage[T1]{fontenc}
\usepackage{amsmath,amssymb}
\usepackage{graphicx}
\usepackage{booktabs}
\usepackage{hyperref}
\usepackage[margin=0.75in]{geometry}
\usepackage{natbib}
\usepackage{xcolor}
\usepackage{enumitem}
\usepackage{microtype}
\usepackage{multirow}

\hypersetup{colorlinks=true, linkcolor=blue!60!black, citecolor=blue!60!black, urlcolor=blue!60!black}

\title{Fashion Florence: Fine-Tuning Florence-2 for\\Structured Fashion Attribute Extraction}

\author{
  Anushree Berlia\\
  \texttt{github.com/anushreeberlia/loom}
}

\date{May 2026}

\begin{document}
\maketitle

% ══════════════════════════════════════════════════════════════════════════════
\begin{abstract}
We present \textbf{Fashion Florence}, a Florence-2 vision--language model
fine-tuned with LoRA to extract structured fashion attributes from clothing
images.  Given a single photograph, the model generates a JSON object
containing category, color, material, style tags, and occasion
tags---structured output suitable for direct programmatic consumption by
downstream recommendation and retrieval systems.

Fine-tuning data is derived from the iMaterialist Fashion dataset (228
labels), where we collapse fine-grained annotations into a compact 6-category,
16-color, 19-style schema via rule-based label engineering.  We apply LoRA
($r{=}16$, $\alpha{=}32$) to all decoder linear layers, training for 3 epochs
on 3,688 examples.

On a held-out test set of 461 images, Fashion Florence achieves
\textbf{94.6\%} category accuracy and \textbf{63.0\%} material accuracy,
compared to \textbf{89.3\%} / \textbf{43.3\%} for GPT-4o-mini.
Fashion Florence produces valid JSON in \textbf{99.8\%} of outputs,
matching GPT-4o-mini, while running at 0.77B parameters on a single GPU at
zero marginal inference cost.  Style tag F1 reaches \textbf{0.753}
vs.\ \textbf{0.398} for GPT-4o-mini, an 89\% relative improvement.  When the model returns \texttt{"unknown"} for color, a
FashionCLIP zero-shot fallback resolves it with negligible additional
latency.

The model is deployed as a Hugging Face Space and integrated into Loom, an
open-source outfit recommendation system.

\smallskip\noindent
Model: \url{https://huggingface.co/anushreeberlia/fashion-florence}\\
Code: \url{https://github.com/anushreeberlia/loom}
\end{abstract}

% ══════════════════════════════════════════════════════════════════════════════
\section{Introduction}
\label{sec:intro}

Fashion attribute extraction---automatically identifying the category, color,
material, and style of a garment from its image---is a foundational task for
fashion recommendation, search, and catalog management.  Downstream systems
require structured, schema-valid outputs: an outfit recommendation engine
needs to know that an item is a ``navy wool blazer'' with style tags
``classic'' and ``workwear,'' not a free-text caption that must be parsed.

Prior approaches to fashion attribute recognition have used multi-label CNNs
\citep{liu2016deepfashion, guo2019imaterialist}, which require separate
classifiers per attribute and cannot generalize to new schema fields.
General-purpose multimodal models like GPT-4V/4o can produce structured
output via careful prompting, but are expensive (\$0.01--0.03 per image),
rate-limited, and introduce latency and reliability concerns for production
pipelines.

We present \textbf{Fashion Florence}, which occupies a middle ground: a
vision--language model (0.77B parameters) fine-tuned to emit
constrained, schema-valid JSON from fashion images.  Our contributions:

\begin{enumerate}[leftmargin=*,topsep=2pt,itemsep=1pt]
\item \textbf{Structured generation framing}: We formulate fashion attribute
  extraction as conditional text generation where the target is a compact
  JSON string, allowing a single model to produce all attributes
  simultaneously in a schema the downstream system can directly parse.

\item \textbf{Label engineering from iMaterialist}: We develop a mapping
  layer that collapses 228 fine-grained iMaterialist labels into a compact,
  application-oriented schema (6 categories, 16 colors, 19 style tags, 15
  occasion tags), deriving style and occasion signals from granular category
  names.

\item \textbf{Practical deployment}: The fine-tuned model runs on a single
  GPU with sub-second inference, is deployed as a Hugging Face Space, and
  includes a FashionCLIP zero-shot color fallback for ambiguous cases.
\end{enumerate}

% ══════════════════════════════════════════════════════════════════════════════
\section{Related Work}
\label{sec:related}

\paragraph{Fashion attribute recognition.}
DeepFashion \citep{liu2016deepfashion} introduced large-scale benchmarks for
clothing recognition with attribute annotations.  The iMaterialist Fashion
dataset \citep{guo2019imaterialist} provides 228 fine-grained labels across
multiple attribute groups.  Traditional approaches train multi-label CNNs,
requiring separate classification heads per attribute type.

\paragraph{Vision--language models for fashion.}
FashionVLP \citep{goenka2022fashionvlp} learns joint vision--language
representations for fashion retrieval.  FashionCLIP
\citep{chia2022fashionclip} adapts CLIP \citep{radford2021clip} to fashion
image--text pairs.  These models produce embeddings rather than structured
attributes, making them complementary to our approach.

\paragraph{Structured output from VLMs.}
Florence-2 \citep{xiao2024florence2} supports multiple vision tasks through
a unified sequence-to-sequence interface.  LoRA \citep{hu2022lora} enables
parameter-efficient fine-tuning.  Recent work has explored structured output
generation from VLMs for document understanding and chart parsing; we apply
this pattern to fashion attribute extraction.

\paragraph{LLM-based attribute extraction.}
GPT-4V and GPT-4o-mini can produce structured fashion descriptions via
careful system prompting.  However, API costs (\$0.01--0.03/image), rate
limits, and non-deterministic output formatting make them impractical as the
sole vision backend for high-throughput catalog processing.  Our approach
provides a self-hosted alternative with deterministic schema compliance.

% ══════════════════════════════════════════════════════════════════════════════
\section{Method}
\label{sec:method}

\subsection{Task Formulation}

We frame fashion attribute tagging as conditional text generation.  Given an
image $I$ and a fixed prompt $p$ = \texttt{"Analyze this clothing item image
and return structured fashion tags as JSON."}, the model generates:

\begin{equation}
  \hat{y} = \arg\max_y \; P(y \mid I, p; \theta)
\end{equation}

\noindent where $\hat{y}$ is a compact JSON string with fields:
\texttt{category}, \texttt{primary\_color}, \texttt{material},
\texttt{style\_tags} (list), and \texttt{occasion\_tags} (list).

The target is serialized with \texttt{json.dumps(separators=(",",":"))} so
the model learns single-line, compact JSON without whitespace variation.

\subsection{Base Model}

Fashion Florence builds on Florence-2-large \citep{xiao2024florence2}, a
0.77B parameter vision--language model with a DaViT vision encoder and an
autoregressive text decoder.  Florence-2's unified sequence-to-sequence
interface---where different vision tasks are distinguished by the input
prompt rather than separate heads---makes it well-suited for generating
structured text conditioned on images.

\subsection{Training Data and Label Engineering}
\label{sec:labels}

We derive training data from the iMaterialist Fashion dataset
\citep{guo2019imaterialist}, accessed via \texttt{Marqo/iMaterialist} on the
Hugging Face Hub.  The original dataset provides 228 fine-grained labels
across groups: category (105 types), color (19), material (30), style,
pattern, sleeve, neckline, and gender.

A rule-based mapping layer collapses these into our target schema:

\begin{itemize}[leftmargin=*,topsep=2pt,itemsep=1pt]
\item \textbf{Categories (6)}: top, bottom, dress, layer, shoes, accessory.
  Approximately 105 iMaterialist categories are mapped; items with no valid
  mapping (swimwear, underwear, costumes) are discarded.
\item \textbf{Colors (16)}: black, white, gray, beige, brown, blue, navy,
  green, yellow, orange, red, pink, purple, metallic, multi, unknown.
\item \textbf{Materials}: Direct mapping from iMaterialist material labels
  (cotton, silk, denim, etc.), used as-is.
\item \textbf{Style tags (19)}: Derived from granular category names and
  iMaterialist style labels.  For example, ``Suits \& Blazers'' $\to$
  \{classic, workwear\}; ``Clubbing Dresses'' $\to$ \{sexy, glamorous\};
  ``Bodycon'' style label $\to$ \{sexy, glamorous\}.
\item \textbf{Occasion tags (15)}: Derived similarly from categories.
  ``Athletic Shirts'' $\to$ \{gym, workout\}; ``Dress Shirts'' $\to$
  \{work\}.  Default: \texttt{["everyday"]} when no rules match.
\end{itemize}

From 10,000 sampled rows, 4,610 pass validation (non-null category mapping)
and are split 80/10/10 into 3,688 / 461 / 461 for train / validation / test.

\begin{table}[h]
\centering
\caption{Category distribution in the training set.}
\label{tab:data}
\begin{tabular}{lrr}
\toprule
Category & Count & \% \\
\midrule
Dress     & 1,684 & 36.5 \\
Top       & 1,450 & 31.5 \\
Bottom    &    917 & 19.9 \\
Layer     &    489 & 10.6 \\
Shoes     &     68 &  1.5 \\
Accessory &      2 &  0.04 \\
\midrule
Total     & 4,610 & 100.0 \\
\bottomrule
\end{tabular}
\end{table}

\subsection{Fine-Tuning with LoRA}

We apply Low-Rank Adaptation (LoRA) \citep{hu2022lora} to all linear layers
of the Florence-2 decoder via the PEFT library.

\begin{table}[h]
\centering
\caption{Training configuration.}
\label{tab:hyperparams}
\begin{tabular}{ll}
\toprule
Parameter & Value \\
\midrule
LoRA rank $r$ & 16 \\
LoRA $\alpha$ & 32 \\
LoRA dropout & 0.05 \\
Target modules & all linear \\
Learning rate & $2 \times 10^{-4}$ \\
Weight decay & 0.01 \\
Warmup ratio & 0.05 \\
Batch size (per device) & 2 \\
Gradient accumulation & 8 \\
Effective batch size & 16 \\
Epochs & 3 \\
Max target length & 256 tokens \\
Precision & FP16 (mixed) \\
Checkpoint selection & lowest eval loss \\
\bottomrule
\end{tabular}
\end{table}

The loss is standard causal language modeling cross-entropy with padding
tokens masked.

\subsection{Zero-Shot Color Fallback}

Fashion Florence occasionally returns \texttt{"unknown"} for
\texttt{primary\_color}, particularly for muted or ambiguous tones.  When
this occurs, we fall back to FashionCLIP \citep{chia2022fashionclip}
zero-shot classification: the image embedding is compared against text
embeddings of \texttt{"a photo of a \{color\} piece of clothing"} for each
of the 16 color labels, and the highest-cosine label is selected.

Since FashionCLIP is already loaded in our downstream pipeline for item
embedding, this adds negligible latency ($<$10\,ms).

% ══════════════════════════════════════════════════════════════════════════════
\section{Experiments}
\label{sec:experiments}

We evaluate Fashion Florence against GPT-4o-mini on the held-out test set
(461 images).  GPT-4o-mini receives the same detailed system prompt used in
our production fallback path, specifying the exact output schema with
per-field instructions and examples.  We additionally contextualize against
Gemini 2.0 Flash results from \citet{kumar2025fashion}, who report macro F1
scores on 18 fashion attributes from DeepFashion-MultiModal using zero-shot
prompting.

\paragraph{Baseline scope.}
We do not include a base (unfine-tuned) Florence-2 baseline because the
pretrained model was not trained to produce structured JSON output in our
schema; without fine-tuning, it generates free-form captions that fail JSON
parsing entirely.  The relevant comparison is therefore against the strongest
available alternative that \emph{can} produce structured output---GPT-4o-mini
with schema-constrained prompting.  An ablation isolating the contribution
of LoRA fine-tuning versus zero-shot Florence-2 with structured prompting is
an important direction for future work.

\subsection{Metrics}

\begin{itemize}[leftmargin=*,topsep=2pt,itemsep=1pt]
\item \textbf{JSON validity}: Percentage of outputs that parse as valid JSON
  with all required fields present.
\item \textbf{Category accuracy}: Exact match on the 6-class category field.
\item \textbf{Material accuracy}: Exact match (case-insensitive).
\item \textbf{Style F1}: Per-sample set F1 between predicted and ground-truth
  \texttt{style\_tags}, averaged over all samples with valid JSON output.
\item \textbf{Occasion F1}: Same as style F1 for \texttt{occasion\_tags}.
\end{itemize}

\noindent Color accuracy is omitted from the primary comparison because
\texttt{primary\_color} is handled by a FashionCLIP zero-shot fallback in
production (Section~\ref{sec:labels}).  For transparency: the Florence model
alone achieves 11.6\% color exact-match on the test set, reflecting the
difficulty of color prediction from images where ground-truth labels are
themselves ambiguous (e.g., ``navy'' vs.\ ``blue'').

\subsection{Results}

\begin{table}[h]
\centering
\caption{Fashion attribute extraction results on held-out test set (N=461).}
\label{tab:results}
\begin{tabular}{lccccr}
\toprule
 & JSON & Cat. & Mat. & Style & Occ. \\
Method & Valid & Acc. & Acc. & F1 & F1 \\
\midrule
GPT-4o-mini & 99.8\% & 89.3\% & 43.3\% & 0.398 & 0.553 \\
\textbf{Fashion Florence} & \textbf{99.8\%} & \textbf{94.6\%} & \textbf{63.0\%} & \textbf{0.753} & \textbf{0.591} \\
\midrule
$\Delta$ (absolute) & 0.0 & +5.3 & +19.7 & +0.355 & +0.038 \\
\bottomrule
\end{tabular}
\end{table}

Fashion Florence outperforms GPT-4o-mini on every accuracy metric while
matching its JSON validity (both 99.8\%).  The largest gain is on style F1
(+89\% relative), reflecting the fine-tuned model's superior understanding
of fashion-specific vocabulary.  Category accuracy improves by 5.3
percentage points; material accuracy by 19.7.

\paragraph{External context.}
\citet{kumar2025fashion} evaluate GPT-4o-mini and Gemini 2.0 Flash on 18
fashion attributes from DeepFashion-MultiModal in a zero-shot setting,
reporting macro F1 of 0.433 and 0.568 respectively.  While their task setup
differs from ours (18 attributes vs.\ 5 fields, different dataset), these
results confirm that zero-shot LLMs struggle with fashion attribute
prediction, and that domain-specific fine-tuning (as in Fashion Florence)
provides substantial improvements.

\subsection{Per-Category Breakdown}

\begin{table}[h]
\centering
\caption{Fashion Florence per-category accuracy on the held-out test set.}
\label{tab:percategory}
\begin{tabular}{lrcc}
\toprule
Category & N & Cat. Acc. & Mat. Acc. \\
\midrule
Top       & 161 & 95.7\% & 67.7\% \\
Dress     & 145 & 95.9\% & 55.9\% \\
Bottom    &  91 & 92.3\% & 65.9\% \\
Layer     &  58 & 93.1\% & 63.8\% \\
Shoes     &   5 & 80.0\% & 60.0\% \\
\midrule
Overall   & 460 & 94.6\% & 63.0\% \\
\bottomrule
\end{tabular}
\end{table}

Category accuracy exceeds 92\% for all categories with $\geq$50 test
samples.  Shoes accuracy (80\%, N=5) is statistically unreliable due to
the extremely small sample size and is included only for completeness;
the 95\% Clopper--Pearson confidence interval is [28.4\%, 99.5\%].
Accessories (N=0 in test) are omitted entirely.

\subsection{Category-then-Rules Ablation}
\label{sec:cat_rules}

To isolate whether the model's style and occasion predictions reflect
genuine visual recognition or merely memorization of category$\to$tag
mappings, we compare the model's direct predictions against a
\emph{category-then-defaults} baseline: take the model's predicted
category, then assign the top-3 most frequent style tags and top-2 occasion
tags for that category (computed from training set statistics).

\begin{table}[h]
\centering
\caption{Style/occasion prediction: model vs.\ category-then-rules baseline.
``Oracle'' uses ground-truth category.  \textit{To be filled after running
\texttt{eval/eval\_category\_then\_rules.py}.}}
\label{tab:cat_rules}
\begin{tabular}{lcc}
\toprule
Method & Style F1 & Occasion F1 \\
\midrule
Fashion Florence (direct) & 0.753 & 0.591 \\
Category $\to$ defaults (predicted cat.) & --- & --- \\
Category $\to$ defaults (oracle cat.) & --- & --- \\
\bottomrule
\end{tabular}
\end{table}

\subsection{Color Pipeline Evaluation}

Table~\ref{tab:color} reports color accuracy across three configurations.
The production pipeline combines Florence with a FashionCLIP zero-shot
fallback for ambiguous cases.

\begin{table}[h]
\centering
\caption{Color prediction accuracy (exact-match and near-synonym match).
\textit{To be filled after running \texttt{eval/eval\_color\_pipeline.py}.}}
\label{tab:color}
\begin{tabular}{lcc}
\toprule
Method & Exact \% & Near \% \\
\midrule
Florence only & 11.6 & --- \\
FashionCLIP zero-shot only & --- & --- \\
Florence + FashionCLIP fallback (production) & --- & --- \\
\bottomrule
\end{tabular}
\end{table}

\subsection{Efficiency}

\begin{table}[h]
\centering
\caption{Inference cost comparison.}
\label{tab:efficiency}
\begin{tabular}{lrrr}
\toprule
Method & Params & Latency & Cost/image \\
\midrule
GPT-4o-mini & --- & 2,457 ms & $\sim$\$0.003 \\
Fashion Florence & 0.77B & $\sim$800 ms & \$0 marginal\textsuperscript{$\dagger$} \\
\bottomrule
\end{tabular}
\end{table}

\noindent\textsuperscript{$\dagger$}\,\$0 marginal cost on an existing GPU
allocation (Hugging Face Spaces T4); fixed GPU hosting costs are amortized
across all inference requests.

Fashion Florence latency of $\sim$800\,ms reflects GPU inference
on the Hugging Face Space (T4).  The 10.5\,s mean observed in our
evaluation includes network round-trip from a local client to the hosted
Space; on-device or co-located inference would reduce this to sub-second.
GPT-4o-mini latency is API round-trip from the same client.

\subsection{Qualitative Analysis}

We observe three common failure patterns:

\paragraph{Material confusion.}
The most frequent errors involve visually similar materials: chiffon
predicted as silk, polyester as satin, and knit as cotton.  These
confusions are understandable given that material identification often
requires tactile rather than visual information.

\paragraph{Multi-style items.}
Items that span multiple style categories (e.g., a structured blazer with
streetwear-influenced details) tend to receive style tags from only one
dominant category.  The model rarely predicts more than 3 style tags, even
when the ground truth contains 4--5.

\paragraph{Category edge cases.}
Jumpsuits and playsuits are occasionally misclassified between ``dress'' and
``bottom,'' reflecting genuine ambiguity in the category taxonomy.  Long
cardigans are sometimes labeled ``dress'' rather than ``layer.''

\noindent Fashion Florence's strongest domain is category prediction for
common categories (top, dress, bottom), where accuracy exceeds 92\%.  The
style F1 advantage over GPT-4o-mini (0.753 vs.\ 0.398) is largest for
fashion-specific vocabulary: GPT-4o-mini tends to default to generic tags
like ``casual'' and ``modern,'' while Fashion Florence more often predicts
domain-specific tags like ``bohemian,'' ``preppy,'' and ``workwear.''

% ══════════════════════════════════════════════════════════════════════════════
\section{Deployment}
\label{sec:deployment}

Fashion Florence is deployed as a Docker container on Hugging Face Spaces,
exposing a \texttt{POST /analyze} endpoint.  The production client
(\texttt{fashion\_florence.py}) handles HF Space cold starts with retry logic
(up to 2 retries, 120s initial timeout) and 503-detection for sleeping
Spaces.

The model outputs all 5 trained fields (\texttt{category},
\texttt{primary\_color}, \texttt{material}, \texttt{style\_tags},
\texttt{occasion\_tags}).  In the Loom production pipeline, the model's
\texttt{occasion\_tags} prediction is overridden by deterministic
style--occasion rules (which achieve higher reliability given correct
category), and three additional fields (\texttt{season\_tags},
\texttt{fit}, \texttt{formality}) are derived via material--season and
style--fit mappings, expanding to the full 8-field Loom schema.

Fashion Florence serves as the default vision backend in Loom, an
open-source outfit recommendation system that uses the extracted attributes
for downstream embedding, retrieval, and outfit scoring.

% ══════════════════════════════════════════════════════════════════════════════
\section{Limitations}
\label{sec:limitations}

\paragraph{Training data imbalance.}
The dataset is heavily skewed toward dresses (36.5\%) and tops (31.5\%),
with shoes (1.5\%) and accessories (0.04\%) severely underrepresented
(Table~\ref{tab:data}).  This reflects iMaterialist's natural distribution
and likely degrades performance on minority categories.

\paragraph{Derived labels and circular evaluation.}
Style and occasion tags are derived from category names via hand-written
rules, not from human annotations of the actual images.  A blazer is always
tagged \{classic, workwear\} regardless of whether the specific image shows a
casual or formal blazer.  Consequently, style F1 and occasion F1 partly
measure the model's ability to reproduce a deterministic category-to-tag
mapping rather than to recognize visual style.  An alternative pipeline
that predicts category first and then applies rules post-hoc might achieve
comparable style/occasion metrics; however, our end-to-end approach
generates all fields jointly, which simplifies the production pipeline and
allows the model to occasionally surface style signals not captured by the
rules (e.g., predicting ``edgy'' for a conventionally categorized top based
on its visual appearance).  Future work should evaluate against
human-annotated style labels to isolate genuine visual style recognition.

\paragraph{Color prediction.}
Color remains the weakest field, particularly for muted, pastel, or
multi-tone garments where the ``correct'' color is ambiguous even to human
annotators.  The FashionCLIP fallback mitigates but does not eliminate this.

% ══════════════════════════════════════════════════════════════════════════════
\section{Conclusion}

We presented Fashion Florence, a Florence-2 model fine-tuned with LoRA to
produce structured JSON fashion attributes from clothing images.  By framing
attribute extraction as conditional text generation with a compact JSON
target, the model provides a self-hosted, sub-second, schema-valid
alternative to general-purpose multimodal APIs for fashion catalog
processing.  The label engineering pipeline from iMaterialist demonstrates
that rich supervision for domain-specific structured generation can be
derived from existing multi-label datasets without additional annotation.

% ══════════════════════════════════════════════════════════════════════════════
\bibliographystyle{plainnat}
\begin{thebibliography}{10}

\bibitem[Chia et~al.(2022)]{chia2022fashionclip}
Chia, P.~J., Attanasio, G., Bianchi, F., Terragni, S., Hung, A.,
Lucaselli, E., Gashteovski, K., Rossiello, G., Subramanian, N., and
Zanzotto, R.
\newblock {FashionCLIP}: Connecting language and images for product
representations.
\newblock \emph{arXiv preprint arXiv:2204.03972}, 2022.

\bibitem[Goenka et~al.(2022)]{goenka2022fashionvlp}
Goenka, S., Zheng, Z., Jain, A., Srikumar, V., Annavaram, M., and Bui, T.
\newblock {FashionVLP}: Vision language transformer for fashion retrieval with
feedback.
\newblock In \emph{CVPR}, 2022.

\bibitem[Guo et~al.(2019)]{guo2019imaterialist}
Guo, S., Huang, W., Zhang, X., Srikhanta, P., Cui, Y., Li, Y., Adam, H.,
Scott, M.~R., and Belongie, S.
\newblock The {iMaterialist} fashion attribute dataset.
\newblock In \emph{CVPR Workshops}, 2019.

\bibitem[Hu et~al.(2022)]{hu2022lora}
Hu, E.~J., Shen, Y., Wallis, P., Allen-Zhu, Z., Li, Y., Wang, S., Wang,
L., and Chen, W.
\newblock {LoRA}: Low-rank adaptation of large language models.
\newblock In \emph{ICLR}, 2022.

\bibitem[Kumar et~al.(2025)]{kumar2025fashion}
Kumar, A., et~al.
\newblock Can {GPT-4o} mini and {Gemini 2.0 Flash} predict fine-grained fashion
product attributes? {A} zero-shot analysis.
\newblock \emph{arXiv preprint arXiv:2507.09950}, 2025.

\bibitem[Liu et~al.(2016)]{liu2016deepfashion}
Liu, Z., Luo, P., Qiu, S., Wang, X., and Tang, X.
\newblock {DeepFashion}: Powering robust clothes recognition and retrieval
with rich annotations.
\newblock In \emph{CVPR}, 2016.

\bibitem[Radford et~al.(2021)]{radford2021clip}
Radford, A., Kim, J.~W., Hallacy, C., Ramesh, A., Goh, G., Agarwal, S.,
Sastry, G., Askell, A., Mishkin, P., Clark, J., Krueger, G., and
Sutskever, I.
\newblock Learning transferable visual models from natural language
supervision.
\newblock In \emph{ICML}, 2021.

\bibitem[Xiao et~al.(2024)]{xiao2024florence2}
Xiao, B., Wu, H., Xu, W., Dai, X., Hu, H., Lu, Y., Zeng, M., Liu, C.,
and Yuan, L.
\newblock {Florence-2}: Advancing a unified representation for a variety of
vision tasks.
\newblock In \emph{CVPR}, 2024.

\end{thebibliography}

\end{document}
